\section{Kill Chain y Value Chain: comparten el esqueleto de software, no las conclusiones éticas}

Sankar considera la kill chain gubernamental y la value chain comercial como decision chains de estructura similar: la organización fija objetivos, construye un operational picture, asigna recursos, ejecuta el workflow y luego observa los resultados. Una plataforma tecnológica puede reutilizar el entity model, el constraint engine, el workflow, el audit y el feedback; sin embargo, la base de autoridad con la que un gobierno emplea la fuerza y la de una empresa para asignar inventario o capital no son la misma, y la semejanza en la estructura del software no justifica borrar las diferencias éticas.

\begin{figure}[H]
\centering
\includegraphics[width=0.94\textwidth]{images/09-common-decision-chain.png}
\caption{Los sistemas gubernamentales y comerciales comparten el esqueleto de sensing, modeling, allocation y feedback, pero conservan authority distintas.}
\figsource{Redibujado a partir de los subtítulos 25:00--30:00, `images/09-common-decision-chain.png`.}
\end{figure}

\begin{knowledgebox}{Lectura de la figura: lo reutilizable son los primitive, no los objetivos finales}
El shared software puede proporcionar ontology, reglas versionadas, resource constraint, workflow y audit trail. Los objetivos gubernamentales están sujetos a la constitución, la ley y la autorización de mando; los objetivos comerciales están sujetos a contratos, regulación y los derechos de empleados y clientes. La plataforma debe hacer que estas reglas sean configurables y auditables, en lugar de trasladar de manera encubierta los objetivos predeterminados de una industria a otra.
\end{knowledgebox}

\begin{table}[H]
\centering
\begin{tabular}{L{0.18\textwidth}L{0.34\textwidth}L{0.38\textwidth}}
\toprule
Primitive & Capacidad general & Lo que debe definirse externamente \\
\midrule
Ontology & Unificar entity, relation y action & Qué clasificaciones son legítimas y quién puede consultarlas \\
Constraint engine & Verificar recursos, políticas y compatibilidad & Origen y prioridad de las reglas \\
Workflow & Coordinar los pasos de personas y sistemas & Qué pasos requieren aprobación humana \\
Optimization & Elegir una opción bajo restricciones & Objective y límites no negociables \\
Audit trail & Registrar entradas, decisiones y ejecución & Plazo de conservación, supervisores y derecho de corrección \\
Feedback & Actualizar la estrategia con los resultados & Qué outcome cuenta como éxito o como daño \\
\bottomrule
\end{tabular}
\caption{La plataforma tecnológica proporciona decision primitive; los objetivos legítimos y los límites de los derechos deben definirlos la organización y las instituciones.}
\end{table}

\subsection{Resumen de la sección}

El gobierno y el sector comercial pueden compartir infraestructura de decisión, pero no objective que no hayan pasado por revisión. Ontology, workflow, constraint y audit son primitive reutilizables; authority, rights y harm threshold deben provenir de las instituciones de cada ámbito. Esta distinción también se aplica a la IA: la oferta de modelos puede ser general, pero la generación de valor depende de un sistema de demanda concreto.
